\hypertarget{sec:objetivos}{}%
\begin{center}
    {\Large\bfseries 3. OBJETIVOS}
\end{center}
\vspace{1em}

\noindent{\large\bfseries 3.1 OBJETIVO GENERAL}\par\vspace{0.5em}
Analizar de qué manera la disponibilidad y el estado operativo de las aplicaciones financieras digitales Nequi, DaviPlata y Bancolombia condicionan la ejecución de las actividades cotidianas e inciden en el grado de dependencia de los estudiantes universitarios de la UPTC, Sede Tunja, mediante la aplicación de un estudio, para identificar las principales afectaciones ante fallas del servicio y reconocer los mecanismos de contingencia implementados.

\vspace{1em}
\noindent{\large\bfseries 3.2 OBJETIVOS ESPECÍFICOS}\par\vspace{0.5em}
\begin{itemize}
    \item[\textbf{a)}] \textbf{Identificar} el nivel de dependencia de los estudiantes respecto a Nequi, DaviPlata y Bancolombia, determinando las actividades cotidianas específicas que resultan más afectadas ante una caída o falla del servicio.\\
    \textit{Meta:} Aplicación de un instrumento de recolección de datos y elaboración de una matriz de caracterización de frecuencia de uso, hábitos de pago y nivel de afectación operacional en los estudiantes de la UPTC, Sede Tunja.

    \item[\textbf{b)}] \textbf{Comparar} Comparar cuál de las tres plataformas analizadas (Nequi, DaviPlata y Bancolombia) presenta un mayor grado de vitalidad e impacto operacional en las actividades diarias de la comunidad estudiantil evaluada.\\
    \textit{Meta:} Matriz comparativa y reporte cuantitativo de vulnerabilidad que clasifique e identifique la aplicación de mayor criticidad para el desarrollo de la rutina universitaria.

    \item[\textbf{c)}] \textbf{Reconocer} los mecanismos de contingencia y las alternativas (uso de efectivo de reserva, migración a otras plataformas o cancelación de pagos) que implementan los estudiantes ante la indisponibilidad de su billetera digital preferida.\\
    \textit{Meta:} Fichas de caracterización de respuestas de contingencia y mapa de hábitos de respaldo frente a la falla de sistemas financieros móviles en el campus universitario.

\end{itemize}